% section | subsection | subsubsection | paragraph | subparagraph | verbatim | enumerate | item

% Just let it be an Article bro
\documentclass{article}

% Cover Metadata
\title{Matrix Algebra for Engineers}
\date{03 Dec 2025}
\author{Maulana Hafidz Ismail}

% Equation Setup
\usepackage{amsmath}

% Image Setup
\usepackage{graphicx}
\usepackage{float}

% Link Setup
\usepackage{hyperref}

% Highlighting Setup
\usepackage{soul}
\usepackage{xcolor}
\sethlcolor{red!30}
\newcommand{\hlblue}[1]{\sethlcolor{cyan!30}\hl{#1}\sethlcolor{yellow}}

% Line Width
\sloppy
\setlength{\emergencystretch}{3em}
\hbadness=99999

% Code Syntax Setup
\usepackage{listings}
\usepackage{xcolor} % untuk warna
\lstset{
basicstyle=\ttfamily\small,
frame=single,
breaklines=true,
backgroundcolor=\color{gray!10},
keywordstyle=\color{blue},
commentstyle=\color{gray},
stringstyle=\color{red},
tabsize=2,
captionpos=b
}

% Start of the Article
\begin{document}

% Cover Page
\pagenumbering{gobble}
\maketitle
\newpage

% Table of Content Page
\tableofcontents
\newpage
\pagenumbering{arabic}

% Main Content
\section{Matrices}
\subsection{Definisi Matriks}
Matrix adalah sebuah kumpulan angka yang disusun dalam baris dan kolom. Digunakan untuk merepresentasikan data atau sistem persamaan. Matriks dinyatakan dalam format m x n, di mana m adalah jumlah baris dan n adalah jumlah kolom.

\subsection{Jenis Matriks Berdasarkan Panjangnya}
\begin{itemize}
    \item Square Matrix: Matriks yang memiliki jumlah baris dan kolom yang sama. Contoh: Matriks 2 x 2.

          \[ A = \begin{bmatrix} a & b \\ c & d \end{bmatrix} \]

    \item Rectangular Matrix: Matriks yang memiliki jumlah baris dan kolom yang berbeda. Contoh: Matriks 3 x 2.

          \[ B = \begin{bmatrix} a & b \\ c & d \\ e & f \end{bmatrix} \]
\end{itemize}

\subsection{Jenis Matriks Berdasarkan Isinya}
\begin{itemize}
    \item Zero Matrix: Matriks yang semua elemennya adalah nol.
    \item Identity Matrix: Matriks persegi yang memiliki 1 di diagonal utama dan 0 di tempat lainnya. Berfungsi sebagai elemen netral dalam perkalian matriks.
    \item Lower Triangular Matrix: Matriks di mana semua elemen di atas diagonal utama adalah nol.
    \item Upper Triangular Matrix: Matriks di mana semua elemen di bawah diagonal utama adalah nol.
    \item Diagonal Matrix: Matriks di mana semua elemen di luar diagonal utama adalah nol.
\end{itemize}

\subsection{Vector (Matrix Satu Baris/Kolom)}
\begin{itemize}
    \item Column Vector: Matriks dengan satu kolom. Notasi umum adalah m x 1.

          \[ x = \begin{bmatrix} a \\ b \\ c \end{bmatrix} \]
    \item Row Vector: Matriks dengan satu baris. Notasi umum adalah 1 x n.

          \[ y = \begin{bmatrix} a & b & c \end{bmatrix} \]
\end{itemize}

\subsection{Notasi Umum untuk Matriks}
Element of a Matrix adalah elemen dalam matriks yang biasanya dinyatakan dengan huruf kecil dan disubscripkan untuk menunjukkan posisi dalam matriks.

\[a_{ij}\]

menunjukkan elemen pada baris i dan kolom j.

Contoh:
\begin{itemize}
    \item Elemen di baris pertama dan kolom pertama dilambangkan sebagai ( \(a_{11}\) ).
    \item Elemen di baris kedua dan kolom kedua dilambangkan sebagai ( \(a_{22}\) ).
\end{itemize}

\subsection{Transpos dan Invers Matriks}
\begin{itemize}
    \item Transpose: Matriks yang diperoleh dengan menukar baris menjadi kolom dan sebaliknya.
    \item Inverse: Matriks yang jika dikalikan dengan matriks asal, menghasilkan Identity Matrix. Hanya ada untuk matriks persegi yang tidak singular.
\end{itemize}

\subsection{Matriks Ortogonal}
Orthogonal Matrix adalah Matriks yang kolom-kolomnya saling tegak lurus dan memiliki panjang satu. Jika dikalikan dengan transposenya, hasilnya adalah Identity Matrix.

\subsection{}
\begin{itemize}
    \item
\end{itemize}

\newpage

\section{Systems of Linear Equations}

\newpage

\section{Vector Spaces}

\newpage

\section{Eigenvalues and Eigenvectors}


% End of Main Content
\end{document}